\documentclass[runningheads]{llncs}
\usepackage{amsmath,amssymb}
\usepackage{graphicx}
\usepackage{booktabs}
\usepackage[colorlinks=false,pdfborder={0 0 0}]{hyperref}

\begin{document}

\title{From Prediction to Decision: Uncertainty-Aware Deferral \\ for Reliable Medical Image Segmentation}

\author{Author Names Withheld}
\institute{Institution Withheld}

\maketitle

\begin{abstract}
Segmentation models produce predictions for every pixel, but clinical deployment
requires knowing which predictions to trust. We formalize this as a selective
prediction problem and compare how uncertainty estimation methods (MC Dropout,
Test-Time Augmentation) and deferral strategies (global thresholding,
image-adaptive, confidence-aware) interact when the goal is reliable decisions
rather than improved segmentation. On retinal vessel analysis (DRIVE, STARE,
CHASE\_DB1), TTA achieves uncertainty AUROC 0.881 versus 0.768 for MC Dropout,
while running 6.3$\times$ faster. Confidence-aware deferral reduces error by
55\% while deferring 12\% of pixels, compared to 26\% error reduction at 33\%
deferral for global thresholding. Temperature scaling does not improve deferral
performance despite marginally reducing calibration error. Cross-dataset
evaluation shows uncertainty quality transfers robustly under domain shift
(Unc-AUROC $>$ 0.83 on external datasets despite 5\% Dice drop).
\end{abstract}

\section{Introduction}

A segmentation model averaging 0.78 Dice across a test set tells a clinician
little about any individual image. What a clinical deployment needs is a policy:
for each pixel, should the output be trusted or deferred to human review?

This question connects uncertainty estimation with selective prediction, two
areas that have developed separately in medical imaging. We address the gap by
studying the \emph{decision layer} on top of existing uncertainty methods.
Rather than proposing new architectures, we ask: given an uncertainty map and a
prediction, what policy should determine which outputs to accept?

Our contributions: (1) A comparison of MC Dropout and TTA evaluated by deferral
effectiveness rather than calibration alone, showing TTA produces substantially
better uncertainty for decisions (Unc-AUROC 0.881 vs.\ 0.768). (2) A
confidence-aware deferral strategy achieving 55\% error reduction at 12\%
deferral rate. (3) Evidence that temperature scaling does not improve
decision quality. (4) Cross-dataset evaluation showing uncertainty quality
transfers more robustly than segmentation accuracy.

\section{Methodology}

\paragraph{Problem setting.} Given input $\mathbf{x}$, a segmentation model
produces prediction $\hat{\mathbf{p}} \in [0,1]^{H \times W}$ and uncertainty
map $\mathbf{u}$. A deferral policy $\delta$ maps these to a binary accept/defer
decision per pixel.

\paragraph{Uncertainty estimation.} MC Dropout ($T = 30$ passes, $p = 0.3$)
uses mutual information: $\text{MI}_{ij} = H(\bar{p}_{ij}) - \frac{1}{T}
\sum_t H(\hat{p}_{ij}^{(t)})$. TTA ($K = 6$ geometric transforms) uses
prediction variance across augmented outputs.

\paragraph{Deferral strategies.} \emph{Global:} defer if $u_{ij} > \tau$.
\emph{Adaptive:} per-image percentile threshold. \emph{Confidence-aware:}
score $s_{ij} = u_{ij} \cdot (1 - c_{ij})$ where $c_{ij} = 2|\hat{p}_{ij}
- 0.5|$, deferring pixels that are both uncertain and near the decision
boundary.

\paragraph{Calibration.} Temperature scaling: $\hat{p}^{\text{cal}} =
\sigma(\text{logit}(\hat{p}) / T)$, with $T$ fit on validation set.

\section{Experiments}

\paragraph{Setup.} U-Net (ResNet-34 encoder) trained on DRIVE (20 train images,
256$\times$256 patches, 80 epochs, Adam, lr $10^{-4}$, Dice+BCE loss).
Cross-dataset: STARE (20 images), CHASE\_DB1 (28 images), zero-shot.

\paragraph{TTA outperforms MC Dropout for deferral.} TTA: Unc-AUROC 0.881,
AUCC (Dice) 0.846, runtime 0.13\,s. MC Dropout: Unc-AUROC 0.768, AUCC 0.801,
runtime 0.82\,s. Segmentation quality is similar (Dice 0.768 vs.\ 0.764).

\paragraph{Confidence-aware deferral is efficient.} At 12\% deferral rate,
confidence-aware reduces error by 55\% (TTA). Global thresholding needs 33\%
deferral for 26\% error reduction (MC Dropout). The confidence-aware score
filters uncertain-but-correct pixels that waste review capacity.

\paragraph{Calibration does not help deferral.} Temperature scaling ($T = 1.19$
for MC, $T = 1.35$ for TTA) barely changes ECE ($<$ 0.001 improvement) and
reduces MC Dropout Unc-AUROC from 0.768 to 0.749.

\paragraph{Uncertainty transfers under domain shift.} DRIVE $\to$ STARE: Dice
drops 5\% (0.764 $\to$ 0.727) but Unc-AUROC rises from 0.768 to 0.846.

\begin{table}[t]
\centering
\caption{Main results on DRIVE. ERR = error reduction ratio.}
\small
\begin{tabular}{lcccccc}
\toprule
Method & Dice & Unc-AUROC & ERR & Def.\ \% & Time (s) \\
\midrule
MC Drop.\ (global) & 0.763 & 0.722 & 26.4\% & 33.4\% & 0.82 \\
MC Drop.\ (conf-aware) & 0.763 & 0.722 & 48.9\% & 10.9\% & 2.65 \\
TTA (adaptive) & 0.768 & 0.881 & 79.5\% & 25.0\% & 0.13 \\
TTA (conf-aware) & 0.768 & 0.881 & 55.2\% & 12.1\% & 0.86 \\
\bottomrule
\end{tabular}
\label{tab:main}
\end{table}

\section{Discussion and conclusion}

TTA's advantage comes from probing input-level prediction fragility rather
than weight-space variance. Confidence-aware deferral succeeds by
distinguishing uncertain-and-borderline pixels from uncertain-but-confident
ones. Temperature scaling fails because it optimizes a distributional property
(ECE) rather than the discriminative property (uncertainty-error correlation)
that deferral requires.

These results argue for evaluating uncertainty methods by their downstream
decision consequences, not calibration metrics alone. Future work should
extend to multi-class segmentation and incorporate human-in-the-loop
evaluation.

\bibliographystyle{splncs04}
\bibliography{../references}

\end{document}
